\documentclass[11pt,a4paper]{article}
% Common preamble for the hint sheets (\input from each file)
\usepackage[utf8]{inputenc}
\usepackage[T1]{fontenc}
\usepackage{lmodern}
\usepackage{amsmath,amssymb,amsthm,mathtools}
\usepackage[a4paper,margin=2.3cm]{geometry}
\usepackage{enumitem}
\usepackage{xcolor}
\usepackage[hidelinks]{hyperref}
\renewcommand{\leq}{\leqslant}
\renewcommand{\le}{\leqslant}
\renewcommand{\geq}{\geqslant}
\renewcommand{\ge}{\geqslant}
\setlength{\parindent}{0pt}
\setlength{\parskip}{0.35em}
% \hintstitle{Sheet title}{Class work / Homework}
\newcommand{\hintstitle}[2]{\begin{center}{\LARGE\bfseries Hints: #1\par}\vspace{0.5em}{\large #2\par}\end{center}
\noindent\textit{The hints to each problem are cumulative. Read one, try again, and open the next only if you are still stuck.}\par\medskip\hrule\medskip}
\newcommand{\problem}[1]{\par\bigskip{\large\bfseries Problem #1}\par\nopagebreak\smallskip}
\newcommand{\subpart}[1]{\par\smallskip\textbf{(#1)}\par\nopagebreak}
\newcommand{\hint}[1]{\par\smallskip\noindent\textcolor{black!70}{\textbf{Hint #1.}}\enspace}
\newcommand{\note}[1]{\par\smallskip\noindent\textit{Note.} #1\par}
\providecommand{\lt}{<}
\providecommand{\gt}{>}
\providecommand{\tightlist}{\setlength{\itemsep}{0pt}\setlength{\parskip}{0pt}}
\emergencystretch=3em

\begin{document}
\hintstitle{Linear algebra -- 10. Matrix norms and localisation of the spectrum}{Homework}

\problem{1}
\hint{1}
Let \(Ax=\lambda x\) and choose \(i\) with \(|x_i|\) maximal. The \(i\)-th equation gives \(|\lambda-a_{ii}|\,|x_i|\le\sum_{j\ne i}|a_{ij}|\,|x_j|\).
\hint{2}
Components: consider \(A(t)=D+t(A-D)\), \(t\in[0,1]\), with \(D\) the diagonal part. The discs grow continuously with \(t\), and the eigenvalues depend continuously on \(t\). An eigenvalue cannot jump from one component to a disjoint one, so the count in a component stays as at \(t=0\). Strict diagonal dominance means \(0\) lies in no disc.
\hint{3}
For the example, take \(A=\begin{pmatrix}0&4\\1&0\end{pmatrix}\) with eigenvalues \(\pm2\). Here row disc \(i\) intersected with column disc \(i\) is \(\{|z|\le1\}\) for both \(i\).

\problem{2}
\hint{1}
Courant--Fischer: if \(\dim L=k\), then \(L\) meets the span of eigenvectors for \(\lambda_k,\dots,\lambda_n\) (dimension \(n-k+1\)). On that span the Rayleigh quotient is \(\le\lambda_k\). Equality is attained on the span of the first \(k\) eigenvectors. Also note the dual form \(\lambda_k=\min_{\dim M=n-k+1}\max_{x\in M}\).
\hint{2}
Cauchy: a principal \((n-1)\times(n-1)\) submatrix is the compression to a coordinate hyperplane \(H\). Taking subspaces \(L\subseteq H\) in the max-min (and in the min-max) gives \(\lambda_k(A)\ge\mu_k\ge\lambda_{k+1}(A)\).
\hint{3}
Weyl: use the min-max form with \(M=M_A\cap M_B\). Here \(M_A\) is spanned by eigenvectors of \(A\) for \(\lambda_i,\dots,\lambda_n\), and \(M_B\) by eigenvectors of \(B\) for \(\lambda_j,\dots,\lambda_n\). Then \(\dim M\ge n-(i+j-1)+1\), and on \(M\) the Rayleigh quotient of \(A+B\) is at most \(\lambda_i(A)+\lambda_j(B)\).

\problem{3}
\hint{1}
Order \(\alpha_1\ge\alpha_2\), \(\beta_1\ge\beta_2\), \(\gamma_1\ge\gamma_2\). Necessary conditions: \(\gamma_1+\gamma_2=\alpha_1+\alpha_2+\beta_1+\beta_2\) (trace), \(\gamma_1\le\alpha_1+\beta_1\), and by Weyl \(\gamma_1\ge\max(\alpha_1+\beta_2,\alpha_2+\beta_1)\).
\hint{2}
Sufficiency: take \(A=\operatorname{diag}(\alpha_1,\alpha_2)\) and \(B=R_\theta\operatorname{diag}(\beta_1,\beta_2)R_\theta^T\) with \(R_\theta\) a rotation. At \(\theta=0\) and \(\theta=\pi/2\) you get the two ends of the interval; by continuity every \(\gamma_1\) in between occurs. Answer: exactly the pairs with the right sum and \(\gamma_1\in[\max(\alpha_1+\beta_2,\alpha_2+\beta_1),\,\alpha_1+\beta_1]\).

\problem{4}
\hint{1}
\(M\) is real symmetric, so all eigenvalues are real. Delete the first row and column: the remaining \(9\times9\) matrix is the standard \(\operatorname{tridiag}(-1,2,-1)\). Its eigenvalues are \(2-2\cos\frac{k\pi}{10}>0\).
\hint{2}
By Cauchy interlacing, \(\lambda_1,\dots,\lambda_9\) of \(M\) are at least the smallest eigenvalue of the submatrix, so they are positive. Finally \(\lambda_{10}\le e_1^TMe_1=-1<0\).

\problem{5}
\subpart{a}
\hint{1}
In terms of the real eigenvalues: \(s=\sum\lambda_i^k=\|\lambda\|_k^k\) because \(k\) is even, and \(t=\sum\lambda_i^{k+1}\) satisfies \(|t|\le\|\lambda\|_{k+1}^{k+1}\). Compare with the monotonicity \(\|\lambda\|_{k+1}\le\|\lambda\|_k\).
\hint{2}
The hypothesis gives \(|t|=s^{(k+1)/k}\), so \(\|\lambda\|_{k+1}=\|\lambda\|_k\). Equality forces at most one non-zero eigenvalue. Then the characteristic polynomial is \(x^{n-1}(x-\operatorname{tr}A)\); apply Cayley--Hamilton.
\subpart{b}
\hint{1}
No. For \(k=1\) the condition is \((\sum\lambda_i)^2=\sum\lambda_i^2\). Try \(A=\operatorname{diag}(2,2,-1)\) and compare \(A^3\) with \(\operatorname{tr}(A)A^2\).

\problem{6}
\hint{1}
Triangularise \(A,B\) simultaneously. The eigenvalues of \(AB\) are \(\alpha_i\beta_i\), all on the unit circle.
\hint{2}
A sum of \(2018\) unit complex numbers equals \(2018\) only if all of them are \(1\). So \(\beta_i=\bar\alpha_i\) and \(\operatorname{tr}B=\overline{\operatorname{tr}A}=\operatorname{tr}A\), since \(A\) is real.

\problem{7}
\hint{1}
\(A+jE\) is singular exactly when \(-j\) is an eigenvalue. How many distinct eigenvalues can \(A\) have?
\hint{2}
\(k=n+1\) works. \(A=\operatorname{diag}(0,-1,\dots,-(n-1))\) shows that \(k=n\) fails. Answer: \(n+1\).

\problem{8}
\hint{1}
\(B=E-A\) commutes with \(A\), so the identity holds for every eigenvalue pair \(a\), \(b=1-a\) (triangularise). Express everything through \(p=ab\) using power sums with \(a+b=1\).
\hint{2}
You get \((1-2p)(1-4p+2p^2)=1-5p+5p^2\), i.e.\ \(p(4p^2-5p+1)=0\). Invertibility excludes \(p=0\), so every \(a_ib_i\in\{1,\tfrac14\}\).
\hint{3}
All combinations occur (diagonal \(A\) with entries \(\frac12\) or \(\frac{1\pm i\sqrt3}2\)). Answer: \(\det(AB)\in\{4^{-m}: 0\le m\le n\}\).

\problem{9}
\hint{1}
By homogeneity, \(A(H_n)=H_n\) means \(\|Ax\|_1=\|x\|_1\) for all \(x\). So \(A\) is an injective linear isometry, and it maps the unit ball \(\operatorname{conv}H_n\) bijectively onto itself.
\hint{2}
A linear bijection of a polytope onto itself maps vertices to vertices. Hence \(Ae_i\in\{\pm e_j\}\), so \(A\) is a signed permutation matrix. There are \(2^nn!\) of these.

\problem{10}
\hint{1}
Write \(A=UD_aU^*\), \(B=VD_bV^*\) and \(W=U^*V\). Then \(\|A-B\|_F^2=\|D_aW-WD_b\|_F^2=\sum_{i,j}|a_i-b_j|^2|w_{ij}|^2\).
\hint{2}
The matrix \((|w_{ij}|^2)\) is doubly stochastic. A linear function on the Birkhoff polytope attains its minimum at a vertex, i.e.\ at a permutation matrix (Birkhoff--von Neumann).
\hint{3}
For the perturbation bound: \(A\,\delta x=\delta b\) gives \(\|\delta x\|\le\|A^{-1}\|\,\|\delta b\|\), and \(\|b\|\le\|A\|\,\|x\|\). Multiply the two inequalities.
\end{document}
